\section{Datos generales del proyecto}

El proyecto \textbf{Sistema Integral de Gestión de Capacitaciones (SIGC-CUSCO)} es una plataforma web creada para registrar, controlar y certificar cursos, talleres y capacitaciones de cualquier institución o empresa. Su objetivo es reemplazar el desorden de las hojas de Excel y listas en papel por un sistema centralizado, rápido y confiable.

\subsection{Objetivo general y específicos}

\paragraph{Objetivo general:}
Desarrollar una aplicación web para administrar todo el ciclo de las capacitaciones (inscripción, control de asistencia por QR, registro de notas y entrega de certificados digitales con validación pública).

\paragraph{Objetivos específicos:}
\begin{itemize}[nosep]
    \item Permitir la inscripción en línea validando el DNI y controlando las vacantes disponibles.
    \item Tomar asistencia de forma rápida en cada clase mediante el escaneo de códigos QR.
    \item Emitir certificados digitales en PDF con código QR para que cualquiera pueda verificar su validez en la web.
    \item Generar actas de notas y reportes claros para la administración y auditoría de los cursos.
\end{itemize}

\subsection{Alcance del sistema}

\paragraph{Lo que sí incluye el sistema:}
\begin{itemize}[nosep]
    \item Inicio de sesión con 3 roles: Administrador, Docente y Participante.
    \item Creación y edición de cursos (fechas, horas, cupos y docente a cargo).
    \item Matrícula de participantes con DNI.
    \item Marcado de asistencia con código QR por sesión.
    \item Registro de calificaciones y descarga de certificados en PDF.
    \item Portal web público para comprobar si un certificado es auténtico.
\end{itemize}

\paragraph{Lo que no incluye (Fuera de alcance):}
\begin{itemize}[nosep]
    \item No incluye pagos con tarjeta en línea (las matrículas se gestionan de forma gratuita o por caja de la institución).
    \item No es un aula virtual para ver videos o foros (tipo Moodle).
    \item No realiza pagos de sueldos a los docentes.
\end{itemize}

\subsection{Justificación del proyecto}

Actualmente, muchas organizaciones manejan sus cursos con listas impresas y diplomas hechos a mano en Word. Esto provoca errores al escribir nombres, pérdida de registros de asistencia y diplomas fáciles de falsificar. Con SIGC-CUSCO, toda la información se guarda en una sola base de datos y cualquier persona o empresa puede escanear el QR de un certificado para saber si es legítimo en pocos segundos.
